% !TeX spellcheck = en_US
\addsection{Deck-building}{\skills/wisdom.png}

\iftoggle{printable}{\vspace{-\baselineskip}}{}

\begin{multicols*}{2}

\subsection*{\pagetarget{Playerdecks}{Player Decks}}
All players have a unique Deck which represents their Main Hero's Abilities and Equipment.
Decks may contain Statistic, Ability, Spell, Artifact and the Main Hero's Specialty Cards.
Each player's Deck starts with 9 cards, built during the game's setup.

\subsection*{General Card Rules}
\begin{enumerate}
  \item Cards can be played \textbf{only on your Turn}, or in a \pagelink{Combat}{Combat} involving your \textbf{Main Hero}.
  \item After a Card is used, discard it.
    Each player has their own separate Discard Pile\index{Discard Pile}.
  \item If Your Deck is empty when you need to draw a Card, \textbf{shuffle your Discard Pile} into a new Deck to draw from.
  \item Whenever your Hero gains a Card for any reason, put it \textbf{directly into your hand} unless otherwise stated.
  \item Whenever you are instructed to \textbf{Search} (X)\index{Search (X)} the Ability, Artifact, or Spell Deck, you may either look at the top (X) Cards from the specified Deck, take one of them to your hand, and discard the others, \textbf{OR} instead of looking at the top (X) Cards, gain the top Card from that Deck's Discard Pile.
  \item The Ability, Artifact, and Spell Decks each have their own Discard Piles, created during the setup, which help you identify these Decks.
    If a Deck ever runs out of Cards, reshuffle it and discard its top Card to form a new Discard Pile.
    Whenever one of these Discard Piles is empty, \textbf{refill it} with that Deck's top Card.
  \item Cards have the following \pagetarget{Card Effects}{types of effects}:
  \begin{itemize}
    \item \textbf{Instant} \svg{instant} Effects are resolved immediately.
    \item \textbf{Activation} \svg{activation} Effects must be played when Activating your own Unit in Combat.
    \item \textbf{Map} \svg{map} Effects cannot be used during Combat.
    \item \textbf{Ongoing} \svg{ongoing} Effects\index{Ongoing Effects} last until they are used up or until the player who played them starts their next Turn (whichever happens first).
    \item \textbf{Permanent} \svg{permanent} Cards\index{Permanent Cards} stay in play until discarded or replaced.
      \textbf{You may only have one permanent Card at a time}; playing another discards the first.
  \end{itemize}
  \item Whenever you are instructed to \textbf{remove a Card}, remove it completely from the game.
\end{enumerate}

\vfill
\hfill{\transparent{0.3}\includegraphics[width=0.9\linewidth]{\art/coin.png}}

\clearpage

\subsection*{\pagetarget{Ability}{Ability} and \pagetarget{Statistic}{Statistic} Cards}

All Ability\index{Ability Card} and Statistic cards\index{Statistic Cards} have a \textbf{Basic Effect}\index{Basic Effect} and a stronger \textbf{Expert} \svgeven{expert} \textbf{Effect}\index{Expert Effect}, which is shown below the Basic Effect.
Whenever you play an Ability or Statistic card, you must choose which effect you are using.
The number of \svgeven{expert} Effects you can use each Round is limited by your Main Hero's \pagelink{Level}{Level}.
Track the number of uses you have in any suitable manner, such as by moving Black Cubes on and off your Hero Card.\par
\bigskip

\vspace{-2em}
\begin{multicols*}{2}
  \footnotesize
  \begin{center}
    \begin{scriptsize}
    \begin{tikzpicture}
      \draw (0, 0) node {\includegraphics[width=\linewidth]{\cards/necromancy_card.png}};
      \draw (0.75, 0.4) node {\encircle{1}};
      \draw (-1, -0.5) node {\encircle{2}};
      \draw (-1, -1.5) node {\encircle{3}};
      \draw (1.5, -0.8) node {\encircle{4}};
      \draw (1.5, -1.9) node {\encircle{4}};
    \end{tikzpicture}
    \end{scriptsize}
    \textbf{\textit{\textcolor{darkcandyapplered}{Ability Card}}}
  \end{center}
  \begin{itemize}
    \item[\textbf{1.}] Name
    \item[\textbf{2.}] Basic Effect
  \end{itemize}
  \columnbreak
  \begin{center}
    \begin{scriptsize}
    \begin{tikzpicture}
      \draw (0, 0) node {\includegraphics[width=\linewidth]{\cards/knowledge.png}};
      \draw (0.75, 0.4) node {\encircle{1}};
      \draw (-1, -0.5) node {\encircle{2}};
      \draw (-1, -1.5) node {\encircle{3}};
    \end{tikzpicture}
    \end{scriptsize}
    \textbf{\textit{\textcolor{darkcandyapplered}{Statistic Card}}}
  \end{center}
  \begin{itemize}
    \item[\textbf{3.}] Expert Effect
    \item[\textbf{4.}] Faction-specific Card*
  \end{itemize}
\end{multicols*}
\bigbreak
\note{11}{
  Certain cards are limited to the Necropolis Faction \svg{necro-note}.
  When a non-Necropolis player draws a Necropolis Card\index{Necropolis Card} from the Ability Deck, they may either discard it and draw a new Card as a replacement or gain it.
  Non-Necropolis players \textbf{cannot use} Faction Specific Cards from their hand in any way besides for effects that discard them.
}

\vspace*{\fill}

\columnbreak

\subsection*{\pagetarget{Artifact}{Artifact} Cards}

Artifact Cards\index{Artifact Card} have a top and bottom effect. When you play one, you must \textbf{choose which effect} you are using.
They are divided into 3 Levels: Minor, Major, and Relic, with different border color.
These Levels relate to the overall Power of the Card and may be referenced when resolving certain effects or during Scenario setup.
Otherwise, all Artifact Cards are normally shuffled together to form the Artifact Deck regardless of their Level.
They are gained through map exploration.\par
Artifacts can be \pagelink{Trading}{traded} in Alliance and Cooperative Scenarios.\par
\begin{multicols*}{2}
  \begin{scriptsize}
  \begin{tikzpicture}
    \draw (0, 0) node {\includegraphics[width=1.2\linewidth]{\cards/artifact.png}};
    \draw (-1.3, 2.1) node {\encircle{1}};
    \draw (1.3, -0.9) node {\encircle{2}};
    \draw (1.4, 2.6) node {\encircle{3}};
    \draw (-1, -1.9) node {\encircle{4}};
  \end{tikzpicture}
  \end{scriptsize}
  \footnotesize
    \null\hfill\textbf{\textit{\textcolor{darkcandyapplered}{Major Artifact}}}
  \columnbreak
  \begin{center}
    \phantom{\ldots}\\
    \phantom{\ldots}
    \begin{scriptsize}
    \begin{tikzpicture}
      \draw (0, 0) node {\includegraphics[width=0.6\linewidth]{\cards/minor_artifact.png}};
      \draw (-0.7, 0.5) node {\encircle{3}};
    \end{tikzpicture}
    \end{scriptsize}\\
    \phantom{\ldots}\textbf{\textit{\textcolor{darkcandyapplered}{Minor Artifact}}}\\
    \vspace{3em}
    \phantom{\ldots}
    \begin{scriptsize}
    \begin{tikzpicture}
      \draw (0, 0) node {\includegraphics[width=0.6\linewidth]{\cards/relic_artifact.png}};
      \draw (-0.7, 0.5) node {\encircle{3}};
    \end{tikzpicture}
    \end{scriptsize}\\
    \phantom{\ldots}\textbf{\textit{\textcolor{darkcandyapplered}{Relic Artifact}}}
  \end{center}
\end{multicols*}
\begin{multicols}{2}
  \footnotesize
  \begin{itemize}
    \item[\textbf{1.}] Name
    \item[\textbf{2.}] Effect
    \columnbreak
    \item[\textbf{3.}] Level
    \item[\textbf{4.}] Fluff
  \end{itemize}
\end{multicols}

\vspace*{\fill}
{\transparent{0.3}\includegraphics[width=\linewidth]{\art/morale_low.png}}

\subsection*{\pagetarget{spells}{Spell Cards}}

Spell Cards\index{Spell Card} have three possible primary effects.
Using the topmost, basic version of the Spell usually has no additional costs.
To access the other effects, you may \textbf{Empower} a Spell\index{Empowering Spells} by paying the indicated cost (3) to get a more powerful outcome (4).
You may pay this cost by playing other cards for their Empower \svg{empower} effect (e.g. the Spell Power Statistic) before casting the Spell.
All Spell Cards also have an alternative bottom (5) \svg{empower} effect.
You cannot preserve the played \svg{empower}: casting a Spell consumes all the \svg{empower} played so far, and if you stop playing cards, the power dissipates.
You may use a cheaper \svg{empower} effect than what you've paid for, though.\par

{
  \centering
  \begin{tikzpicture}
    \draw (0, 0) node {\includegraphics[width=0.6\linewidth]{\cards/spell.png}};
    \draw (-1.3, 2.5) node {\encircle{1}};
    \draw (-1.8, 0.3) node {\encircle{2}};
    \draw (-1.1, -1) node {\encircle{3}};
    \draw (1.1, -1) node {\encircle{4}};
    \draw (0.8, -2.8) node {\encircle{5}};
  \end{tikzpicture}\\
  \medskip
  \footnotesize
  \textbf{\textit{\textcolor{darkcandyapplered}{Spell Card}}}\par

  \noindent
  \begin{minipage}[t]{0.5\linewidth}
    \begin{itemize}[itemsep=0pt, parsep=5pt, topsep=0pt, partopsep=0pt]
      \item[\textbf{1.}] Spell Name
      \item[\textbf{2.}] School of Magic
      \item[\textbf{3.}] Cost to Empower
    \end{itemize}
  \end{minipage}
  \begin{minipage}[t]{0.4\linewidth}
    \begin{itemize}[itemsep=0pt, parsep=5pt, topsep=0pt, partopsep=0pt]
      \item[\textbf{4.}] Spell Effect
      \item[\textbf{5.}] Alternative\\Effect
    \end{itemize}
  \end{minipage}
}
\medskip

Spells can be gained by Building the Mage Guild\index{Mage Guild}.
Building it also enables you to use the Spell Book Token\index{Spell Book Token} for buying additional Spells.
The Token cannot be used during the \textbf{same Round} when the Guild is constructed. Spells can also be \pagelink{Trading}{traded} in Alliance and Cooperative Scenarios.

During Combat, \textbf{only one Spell Card} may be played by each player \textbf{per Combat Round}.

\columnbreak
When you use an effect which lets you cast a Spell from somewhere other than your hand (such as directly from the Spell Deck), place that Spell into your Discard Pile after casting it.

\begin{center}
  \vspace*{\fill}
  {\transparent{0.3}\includegraphics[width=0.9\linewidth]{\art/air_magic.png}}
  \vspace*{\fill}
\end{center}

\subsubsection*{Schools of Magic}
All Spell Cards belong to one School of Magic\index{Schools of Magic}: either Air, Fire, Earth or Water.
Spells with one School symbol on them are called \textbf{Basic Spells}\index{Basic Spell}, while Spells with four identical symbols are called \textbf{Expert Spells}\index{Expert Spell}.
\textbf{Magic Arrow} is a Basic Spell; when casting the \textbf{Magic Arrow}, you must select which School it belongs to.
By default there's no difference, but certain game effects can refer to specific Schools and to Basic/Expert Spells.

\begin{minipage}[t]{0.48\textwidth}
  \centering
  \minipage[b]{0.5\textwidth}
    \centering
    \includegraphics[width=0.75\linewidth]{\images/school_of_fire.png}
    \textit{\textbf{\textcolor{darkcandyapplered}{School of Fire}}}
  \endminipage
  \minipage[b]{0.5\textwidth}
    \centering
    \includegraphics[width=0.75\linewidth]{\images/school_of_water.png}
    \textit{\textbf{\textcolor{darkcandyapplered}{School of Water}}}
  \endminipage
  \hfill\allowbreak%
  \bigbreak
  \minipage[b]{0.5\textwidth}
    \centering
    \includegraphics[width=0.75\linewidth]{\images/school_of_air.png}
    \textit{\textbf{\textcolor{darkcandyapplered}{School of Air}}}
  \endminipage
  \minipage[b]{0.5\textwidth}
    \centering
    \includegraphics[width=0.75\linewidth]{\images/school_of_earth.png}
    \textit{\textbf{\textcolor{darkcandyapplered}{School of Earth}}}
  \endminipage
  \bigbreak
\end{minipage}

\clearpage
\subsection*{\pagetarget{Specialty}{Hero Specialty Cards}}

Hero Specialty Cards\index{Specialty Card} are gained from \pagelink{Level}{Level ups}.
Each Main Hero has a unique set of Specialty Cards.
While many of these Cards have effects which resemble Spell Cards and even make use of \svg{empower} and Schools of Magic, Specialty Cards are their \textbf{own unique category of cards}.
For example, the limit of 1 Spell per Combat Round doesn't apply to Specialty Cards.

{
  \medskip
  \centering
  \includegraphics[width=0.7\linewidth]{\cards/specialty.png}\\
  \medskip
  \footnotesize\textit{A Level 4 Specialty Card, belonging to Catherine the Knight.}
}

\vspace*{\fill}
{\transparent{0.3}\includegraphics[width=\linewidth]{\art/morale_high.png}}
\vspace*{\fill}

\columnbreak
\includegraphics[width=\linewidth]{\art/rust_dragon.jpg}

\clearpage
\subsection*{Searching Example}

\textit{Bob, playing as Deemer the Warlock, has built a Mage Guild during the previous Game Round, enabling him to now use his Spell Book Token to purchase additional Spells.}

\includegraphics[width=\linewidth]{\examples/deemer_receiving_spellbook.png}\par
\textit{He spends the token and pays the cost of 5 \svg{gold} indicated on his Town Board, allowing him to \textbf{Search (2)} the Spell Deck.}\par
\includegraphics[width=\linewidth]{\examples/deemer_spell_cost.png}\par
\filbreak
\textit{He now has the choice of either gaining the top Card of the Spell Card Discard Pile, or drawing the top two Cards of the Spell Deck and gaining one of them.
He is not interested in gaining the Spell Curse, which is on top of the Discard Pile, and instead decides to draw two Cards – a Magic Arrow and a Fireball.}
\includegraphics[width=\linewidth]{\examples/fireball_over_magic_arrow.png}\par
\textit{He decides to keep the Fireball, placing it into his hand and discarding the Magic Arrow into the Spell Discard Pile.}

\vspace*{\fill}

\hfill{\transparent{0.3}\includegraphics[width=0.8\linewidth]{\art/fire_magic.png}}

\end{multicols*}
